\documentclass[10pt,a4paper]{amsart}
\usepackage[margin=19mm]{geometry}
\usepackage{amsmath,amssymb,mathtools}
\usepackage[scr=rsfs,cal=euler]{mathalfa}
\usepackage{xfrac}
\usepackage[unicode,hidelinks,bookmarks=false]{hyperref}
\newcommand{\ie}{i.e.\ }
\newcommand{\cf}{cf.\ }
\newcommand{\resp}{respectively\ }
\newcommand{\Subsection}{\subsection}
\providecommand{\nobreakdash}{\nobreak\hbox{-}}
\setlength{\emergencystretch}{2em}
\multlinegap0pt
\usepackage{amssymb}
\usepackage{mathtools}
\usepackage{dsfont}
\usepackage{enumerate}
\newtheorem{lemma}{Lemma}
\title{Separable neighbourhood of identity in C$^*$-algebras}
\author{Mizanur Rahaman, Mateusz Wasilewski}
\date{Source: arXiv:2603.29556v1 (CC BY 4.0)}
\begin{document}
\maketitle

\noindent\emph{Source and licence.} Separable neighbourhood of identity in C$^*$-algebras. Version arXiv:2603.29556v1 (CC BY 4.0), distributed by arXiv under the Creative Commons Attribution 4.0 International licence, http://creativecommons.org/licenses/by/4.0/, as recorded in the arXiv licence metadata for this exact version and shown on the abstract page https://arxiv.org/abs/2603.29556v1. Full licence notice: \texttt{licenses/arxiv-2603.29556-CC-BY-4.0.txt}.\par\smallskip

\noindent\textit{Editorial scope.} Counted unit: the article's lemma lemma (source0.tex lines 246--248). The article's complete original proof of it is delivered with it (source0.tex lines 249--256, one \texttt{proof} environment of 8 lines). No mathematics is written, repaired or re-derived: the statement and the proof are the article's own lines, in the article's own notation, and no retained line is substituted or re-flowed.  Retained uncounted as the source-local context the proof needs: lines 221--223.  Nothing was omitted from any retained span, so the package carries no omission record.  No retained line contains a citation, so the wrapper carries no bibliography and no bibliography entry.  No retained line contains a figure environment, an image-inclusion command or drawing code, so no image is delivered and the package declares no asset dependency.  Editorial formatting only, in addition: an AMS \texttt{amsart} preamble that restates the article's own declarations and macros used by the retained lines, namely the environments \texttt{lemma} on separate counters, together with the standard packages those lines need.  The retained environments print in this document as Lemma 1 in order of appearance, whereas the article prints the same items with the numbering of its own full document.  The complete native source of the article as distributed in the arXiv e-print for this exact version is bundled unchanged as \texttt{sources/197521/source0.tex}.
\begin{lemma}\label{lem:findimhorodecki}
Let $x \in M_d \otimes N$. Then $x\in Sep_{\sigma}(M_d\otimes N)$ if and only if $(Id\otimes \phi)(x) \in M_d\otimes M_d$ is positive for all normal positive maps $\phi: N \to M_d$. 
\end{lemma}

\begin{lemma}
Let $M$ be a finite, injective, $\sigma$-finite von Neumann algebra. Then $x\in Sep_{\sigma}(M\overline{\otimes} N)$ if and only if $(Id\otimes \phi)(x) \in M\overline{\otimes} M$ is positive for all positive, normal, finite rank maps $\phi: N \to M$.
\end{lemma}

\begin{proof}
Since $M$ is finite and $\sigma$-finite, it admits a faithful normal tracial state $\tau$. It follows that any von Neumann subalgebra $M' \subset M$ admits a normal faithful conditional expectation $\mathbb{E}_{M'}: M \to M'$. As $M$ is injective, it is an increasing union of finite dimensional subalgebras $M_{k}$ with the corresponding conditional expectations $\mathbb{E}_{k}$.

Let now $x \in M \overline{\otimes} N$ satisfy that $(Id\otimes \phi)(x) \in M\overline{\otimes} M$ is positive for all positive, normal, finite rank maps $\phi: N \to M$. Now let's define $x_k:=(\mathbb{E}_{k}\otimes Id)(x)$, then $x_k\rightarrow x$ in the weak$^*$ sense as $\mathbb{E}_{k}$ converges in the point-weak$^{\ast}$ topology to identity. Now note that 
\[(Id\otimes \phi)(x_k)=(Id\otimes \phi)(\mathbb{E}_{k}\otimes Id)(x)=(\mathbb{E}_{k}\otimes Id)(Id\otimes \phi)(x)\geq0,\]
where we have used that $(Id\otimes \phi)(x)\geq 0$ and the fact that $E_k$ is a completely positive map. 
 It follows from Lemma \ref{lem:findimhorodecki} that $x_k$ is separable for all $k \in \mathbb{N}$. So $x$ is separable as well as the set of separable elements is weak$^*$ closed. 
\end{proof}

\end{document}
